\subsection{Using Ordinal Predictions as Surrogate Outcomes in Causal Inference}
\label{sec:surrogate_causal}

A primary motivation for generating high-frequency, building-level wealth trajectories is to evaluate the local impact of clean policy interventions (e.g., new transit infrastructure, rezoning, or place-based policies) in cities lacking yearly survey data. However, using a machine learning prediction as the outcome variable in a causal event study requires defining the conditions under which identification is valid. 

Drawing on the literature for surrogate outcomes \citep{athey2019surrogate}, let $D_i \in \{0,1\}$ denote a policy intervention. Let $W_{i,t}(d)$ be the true latent wealth under treatment state $d$, and $X_{i,t}(d)$ be the observed physical built environment in aerial imagery. The true causal parameter of interest is the Average Treatment Effect on the Treated (ATT) for actual wealth:
\begin{equation}
    \tau_{true} = \mathbb{E}\left[ W_{i,t}(1) - W_{i,t}(0) \mid D_i = 1 \right]
\end{equation}

Because $W_{i,t}$ is unobserved at high frequencies, an econometrician must use the model's conditional expectation $\hat{R}_{i,t}(d) = \mathbb{E}[W_{i,t} \mid X_{i,t}(d)]$, yielding the feasible proxy estimator:
\begin{equation}
    \tau_{proxy} = \mathbb{E}\left[ \hat{R}_{i,t}(1) - \hat{R}_{i,t}(0) \mid D_i = 1 \right]
\end{equation}

By definition, true wealth comprises the capitalized physical signal plus an orthogonal residual: $W_{i,t}(d) = \hat{R}_{i,t}(d) + \epsilon_{i,t}(d)$. Substituting this into the true ATT yields the relationship between the empirical estimate and the true economic effect:
\begin{equation}
    \tau_{true} = \tau_{proxy} + \mathbb{E}\left[ \epsilon_{i,t}(1) - \epsilon_{i,t}(0) \mid D_i = 1 \right]
\end{equation}

For $\tau_{proxy}$ to unbiasedly identify $\tau_{true}$, the second term (the surrogate bias) must equal zero. This condition requires that the policy intervention does not systematically alter local wealth without correspondingly altering the physical built environment. 

As discussed in Section \ref{sec:weakness}, pure demographic sorting into existing housing stock (gentrification without physical construction) is inherently invisible to the visual model. In such cases, $\mathbb{E}[\epsilon_{i,t}(1) - \epsilon_{i,t}(0) \mid D_i = 1] > 0$, meaning the ML-based event study will act as a conservative \textit{lower bound} on total wealth effects. However, if the econometrician specifically defines the estimand as the policy's causal effect on \textit{capitalized structural wealth}---the permanent physical manifestation of economic change---then $\tau_{proxy}$ identifies the exact parameter of interest, making zero-shot causal evaluation highly viable.